\seccionreporte{Implementación del microservicio}{sec:implementacion}

\propositoseccion{%
  Resume la estructura del repositorio y los módulos críticos. No pegue todo el código: cite rutas y fragmentos representativos.%
}

\subsection{Repositorio}

El código fuente completo se encuentra documentado y estructurado en el repositorio oficial:

URL: \url{https://github.com/eduartrob/integratorProjectClustering-back-corvus} (rama \texttt{Palomeque}).

\begin{tabular}{@{}lp{9.5cm}@{}}
  \toprule
  Ruta & Contenido esperado \\
  \midrule
  \texttt{app/} & Código central del servicio, incluyendo enrutamiento, lógica de clustering y NLP. \\
  \texttt{app/models/} & Almacena los artefactos serializados (\texttt{.joblib}) y la base de datos SQLite de auditoría. \\
  \texttt{notebooks/} & Contiene el script \texttt{run\_training.py} y las gráficas exploratorias. \\
  \texttt{collections/} & Colección de Postman estructurada para validar todos los endpoints. \\
  \texttt{Dockerfile} & Configurado para el despliegue directo en el servidor de Contabo. \\
  \texttt{README.md} & Guía detallada de instalación y comandos operativos. \\
  \bottomrule
\end{tabular}

\subsection{Arranque del servicio}

Para inicializar el servicio en el servidor de producción o en un entorno local, el proceso se basa en el estándar de entornos virtuales y FastAPI:

\begin{lstlisting}[style=bashstyle]
# Creacion y activacion del entorno virtual
python3 -m venv venv
source venv/bin/activate

# Instalacion de dependencias
pip install -r requirements.txt
python -m spacy download es_core_news_sm

# Arranque del servidor expuesto en el puerto 3000
uvicorn app.api.server:app --host 0.0.0.0 --port 3000 --reload
\end{lstlisting}

\subsection{Fragmento representativo}

A continuación, presento un fragmento clave del enrutador principal que ilustra nuestra lógica de validación estricta y el registro de auditoría de inferencias:

\begin{lstlisting}[style=pythonstyle]
# app/api/routes.py (extracto)
@router.post("/pre-validate-proposal")
async def pre_validate_proposal(
    user_id: str = Form(...),
    file: UploadFile = File(...)
):
    # Validamos las secciones obligatorias del documento
    section_check = _validate_project_sections(full_text)
    if not section_check["is_valid"]:
        raise HTTPException(
            status_code=400,
            detail=f"Sección obligatoria ausente: {section_check['missing_mandatory'][0]}"
        )
    
    # [Generacion de embeddings y busqueda de similitudes]
    
    # Auditoría: persistencia de la inferencia en SQLite
    log_inference(
        user_id=user_id,
        filename=file.filename,
        score_colision=max_similitud_pct,
        nivel_riesgo=collision_risk_level,
        academic_alignment=academic_alignment
    )
    return quick_analysis
\end{lstlisting}

\subsection{Manejo de errores}

\begin{itemize}
  \item \texttt{400 Bad Request}: Retornado cuando el documento carece de secciones obligatorias (Objetivos, Alcance) o cuando la extracción de texto falla al procesar archivos ilegibles.
  \item \texttt{422 Unprocessable Entity}: Capturado automáticamente por FastAPI cuando la solicitud no incluye los parámetros o archivos requeridos por el esquema.
  \item \texttt{500 Internal Server Error}: Implementamos bloques de captura (\texttt{try-except}) para manejar excepciones imprevistas del modelo de lenguaje o saturación de memoria, garantizando que el cliente reciba un mensaje claro sin comprometer la estabilidad del sistema.
\end{itemize}

\subsection{Dependencias principales}

Para garantizar un rendimiento óptimo y un contenedor ligero, limitamos las dependencias estrictamente a las necesarias:
\begin{itemize}
  \item \texttt{fastapi} y \texttt{uvicorn}: Seleccionadas por su alto rendimiento y manejo asíncrono para el desarrollo de la API.
  \item \texttt{sentence-transformers}: Empleado para la generación de embeddings de 384 dimensiones.
  \item \texttt{umap-learn} y \texttt{hdbscan}: Esenciales para la reducción de dimensionalidad y el clustering basado en densidades.
  \item \texttt{spacy}: Implementado para el reconocimiento de entidades nombradas (NER) y la anonimización de los datos procesados.
  \item \texttt{chromadb}: Utilizado como base de datos vectorial residente en memoria.
\end{itemize}
